\part{Two drivers: heartbeat and ears}

\chapter{The timer}

\begin{goals}
  \item How the PIT works and how to program a frequency
  \item Why a system needs a heartbeat
  \item What a ``tick'' is and what is measured with it
\end{goals}

\section{Why a timer matters so much}

Without a periodic interrupt, an operating system cannot:

\begin{itemize}
  \item take the processor away from a program that will not give it up,
  \item measure time at all,
  \item wake a sleeping process,
  \item time out a device that stopped responding.
\end{itemize}

The timer is what turns a passive kernel into an active one. It is the only thing in the system
that happens \emph{without anyone asking}.

\section{The PIT 8253/8254}

The chip has an input clock of 1,193,182 Hz --- an absurd number that comes from dividing the
original PC's colour television crystal. You cannot choose the frequency directly; you choose a
\term{divisor}, and the chip fires at $1{,}193{,}182 / \text{divisor}$ Hz.

\filetag{lytlnybl/kernel/drivers/timer.c}
\begin{lstlisting}[style=c]
void timer_init(uint32_t frequency) {
    freq = frequency;
    uint16_t divisor = 1193182 / frequency;

    /* tell the PIT how we will send the divisor, and which mode */
    outb(PIT_COMMAND, PIT_ACCESS_LOHIBYTE | PIT_MODE3 | PIT_CHANNEL0 | PIT_BINARY);
    io_wait();

    /* the divisor is 16 bits but the port is 8: send it in two halves */
    outb(PIT_CHANNEL0_DATA, divisor & 0xFF);        /* low byte first */
    io_wait();
    outb(PIT_CHANNEL0_DATA, divisor >> 8);          /* then the high byte */
}
\end{lstlisting}

The project calls \cod{timer\_init(100)}, so the divisor is $1{,}193{,}182 / 100 = 11{,}931$ and the
interrupt fires 100 times per second. \hi{One tick equals 10 milliseconds.}

\begin{warn}{The frequency is not exact}
$1{,}193{,}182 / 11{,}931 = 100.0069$ Hz. The error is tiny but it accumulates: after an hour the
clock has drifted by about a quarter of a second. Systems that need accurate time use a different
source and periodically correct. For a teaching system, and for scheduling, it is irrelevant.
\end{warn}

\begin{concept}{Choosing the frequency is a trade-off}
Higher frequency means finer time resolution and snappier process switching, but more interrupts
per second and therefore more overhead: every tick costs a full save and restore of state. Lower
frequency is cheaper but makes the system feel sluggish and coarse. 100 Hz is a classic
compromise --- it was Linux's default for years.
\end{concept}

\section{What happens on every tick}

\filetag{lytlnybl/kernel/drivers/timer.c}
\begin{lstlisting}[style=c]
volatile uint32_t ticks = 0;

void timer_handler(registers_t* regs) {
    ticks++;

    /* wake any process whose sleep has expired */
    process_t* traversal_process = process_head;
    while (traversal_process) {
        if (traversal_process->state == PROCESS_SLEEPING
            && ticks >= traversal_process->wake_tick) {
            traversal_process->state = PROCESS_READY;
        }
        traversal_process = traversal_process->next;
    }

    schedule(regs);     /* maybe switch process */
}
\end{lstlisting}

Three jobs, in order: advance the clock, wake sleepers, and offer the scheduler a chance to switch.
That last call is where preemptive multitasking actually happens, and Part IX takes it apart.

\begin{concept}{Why \cod{volatile} on \cod{ticks}}
\cod{ticks} is modified inside an interrupt handler and read from ordinary code. Without
\cod{volatile}, a loop like \cod{while (ticks < target);} could be optimised into ``read once, loop
forever'' --- the compiler has no way to know an interrupt changes it. This is the same reason as
the VGA buffer in Chapter 19, from the opposite direction.
\end{concept}

\section{Waiting without a scheduler}

Very early in boot, before processes exist, the kernel sometimes needs to wait. It does so crudely:

\filetag{lytlnybl/kernel/drivers/timer.c}
\begin{lstlisting}[style=c]
void timer_wait_ms(uint32_t ms) {
    uint32_t start = ticks;
    while ((ticks - start) < ms) {
        asm volatile ("hlt");
    }
}
\end{lstlisting}

\cod{hlt} halts the processor until the next interrupt. It is far better than an empty spin loop: it
consumes no power and, in a virtual machine, does not peg a host CPU core at 100\,\%. The timer
interrupt wakes it 100 times a second, the condition is rechecked, and off it goes again.

\begin{warn}{The name lies}
The parameter is called \cod{ms} but is really counted in ticks. With a 100 Hz timer, one tick is
10 ms, so \cod{timer\_wait\_ms(10)} waits 100 ms, not 10. It is a real naming bug in the project.
Harmless today because it is used once during boot, but exactly the kind of thing that wastes an
afternoon later.
\end{warn}

% ============================================================
\begin{recipe}{measuring time with the PIT}{ch23-timer}
Times a piece of work by subtracting two tick readings, then reprograms the PIT to 1000 Hz and
times the same work again. The number of ticks changes by a factor of ten while the work does not:
the ruler changed, not the thing being measured.

It also demonstrates the naming bug above --- \cod{timer\_wait\_ms(100)} is measured taking 100
ticks, that is a full second, not 100 ms.
\end{recipe}

% ============================================================
\chapter{The keyboard}

\begin{goals}
  \item What a scancode is and why it is not ASCII
  \item How press and release are distinguished
  \item How modifier keys are tracked
  \item How a keystroke reaches a blocked user process
\end{goals}

\section{Scancodes are not characters}

When you press a key, the keyboard does not send a letter. It sends a \term{scancode}: a number
identifying \emph{a physical key in a physical position}. The key that says \cod{A} on a UK layout
sends the same scancode as the one that says \cod{Q} on a French layout, because it is the same
physical switch.

Translating scancode to character is entirely the operating system's job, and it is what makes
keyboard layouts possible.

\begin{center}
\begin{tikzpicture}[font=\sffamily\scriptsize,
  b/.style={draw=accent,fill=accentlight,rounded corners=2pt,minimum width=2.5cm,minimum height=0.85cm,align=center}]
  \node[b] (a) {physical key\\{\tiny you press it}};
  \node[b,right=0.8cm of a] (c) {scancode\\{\tiny e.g. 0x1E}};
  \node[b,right=0.8cm of c] (d) {keymap table\\{\tiny lookup}};
  \node[b,right=0.8cm of d] (e) {character\\{\tiny 'a' or 'A'}};
  \foreach \x/\y in {a/c,c/d,d/e} \draw[-{Stealth[length=2.2mm]},accent,thick] (\x)--(\y);
\end{tikzpicture}
\end{center}

\section{Press and release}

Every key generates \emph{two} events: one when pressed (\term{make}) and one when released
(\term{break}). The distinction is encoded in bit 7:

\begin{center}
\begin{tabularx}{0.92\textwidth}{@{}l l X@{}}
\toprule
\textbf{Scancode} & \textbf{Bit 7} & \textbf{Meaning} \\
\midrule
\hexa{1E} & 0 & The \cod{A} key was pressed \\
\hexa{9E} & 1 & The \cod{A} key was released ($\hexa{1E} + \hexa{80}$) \\
\bottomrule
\end{tabularx}
\end{center}

\filetag{lytlnybl/kernel/drivers/keyboard.c}
\begin{lstlisting}[style=c]
if (scancode & KEY_RELEASED) {   /* KEY_RELEASED is 0x80 */
    return;                      /* we ignore releases, except for modifiers */
}
\end{lstlisting}

\section{Modifier keys}

Shift and Control are different: what matters is not the event but the \emph{state}. The driver
keeps a boolean for each and updates it on both press and release:

\filetag{lytlnybl/kernel/drivers/keyboard.c}
\begin{lstlisting}[style=c]
bool keyboard_modifier_keys(uint8_t scancode) {
    if (scancode == SHIFT_PRESS) {
        shift_pressed = true;
        return true;
    } else if (scancode == (SHIFT_PRESS | KEY_RELEASED)) {
        shift_pressed = false;
        return true;
    }

    if (scancode == L_CTRL_PRESS) {
        ctrl_pressed = true;
        return true;
    } else if (scancode == (L_CTRL_PRESS | KEY_RELEASED)) {
        ctrl_pressed = false;
        return true;
    }

    return false;
}
\end{lstlisting}

Returning \cod{true} means ``I handled this, produce no character''. Two keymap tables then decide
the result:

\begin{lstlisting}[style=c]
char c[2];
c[0] = keymap[scancode];
c[1] = '\0';
/* ... */
if (shift_pressed) {
    c[0] = shift_keymap[scancode];
}
\end{lstlisting}

\begin{warn}{Extended scancodes}
Arrow keys, the right-hand Control, Home, End and others send a \hexa{E0} prefix followed by the
real code. The driver detects the prefix and reads a second byte, but currently does nothing with
it:
\begin{lstlisting}[style=c]
void keyboard_extended_scancodes() {
    uint8_t scancode = inb(PS2_DATA);
    if (scancode == 0x5B) { }   /* left GUI key: recognised, ignored */
}
\end{lstlisting}
Implementing arrow keys --- and with them command history and line editing in the shell --- starts
right here. It is one of the most rewarding extensions to attempt.
\end{warn}

\section{From keystroke to blocked process}

This is the most interesting part of the driver, because it is where two subsystems meet. When a
user program calls \cod{read(0, buffer, n)}, the kernel cannot answer immediately --- nobody has
typed anything yet. So it \emph{blocks} the process:

\filetag{lytlnybl/kernel/interrupts.c}
\begin{lstlisting}[style=c]
case SYSCALL_READ:
    /* ... */
    current_process->state = PROCESS_BLOCKED;
    current_process->reading_state.buffer = (void*)regs->ecx;
    current_process->reading_state.size   = regs->edx;
    current_process->reading_state.count  = 0;
    context_switch(current_process, get_next_process(), regs);
    break;
\end{lstlisting}

The process is parked with a note attached: \emph{where} its buffer is, \emph{how big} it is, and
\emph{how much} has been filled so far. Then the kernel switches to another process. The blocked one
is simply never scheduled.

Later, a key arrives and the keyboard handler walks the process list looking for anyone waiting:

\filetag{lytlnybl/kernel/drivers/keyboard.c}
\begin{lstlisting}[style=c]
process_t* traversal_process = process_head;
while (traversal_process) {
  if (traversal_process->state == PROCESS_BLOCKED
      && traversal_process->reading_state.size > 0) {

    void* buffer = traversal_process->reading_state.buffer;
    uint32_t size = traversal_process->reading_state.size;

    /* a rather crude way of doing this (really only the context switcher
     * should be changing cr3), but it works */
    set_cr3((uintptr_t)traversal_process->page_directory);
    ((char*)(buffer))[traversal_process->reading_state.count++] = c[0];
    set_cr3((uintptr_t)kernel_directory);

    if (traversal_process->reading_state.count >= size || c[0] == '\n') {
      traversal_process->regs.eax = traversal_process->reading_state.count;
      traversal_process->reading_state.count = 0;
      traversal_process->reading_state.size = 0;
      traversal_process->state = PROCESS_READY;   /* wake it up */
    }
  }
  traversal_process = traversal_process->next;
}
\end{lstlisting}

\begin{concept}{Why the CR3 dance}
The buffer address (say \hexa{BFFFEF00}) is a \emph{virtual} address, and it only means anything
inside that process's own address space. The keyboard interrupt, however, can arrive while a
completely different process is running --- so \reg{CR3} points somewhere else and that address
means something else entirely, or nothing at all.

So the handler temporarily switches \reg{CR3} to the sleeping process's page directory, writes the
byte, and switches back. It works because the kernel is identity-mapped into every process's
directory, so the code keeps executing correctly throughout.

The comment in the source is honest about this being crude, and it is right. A cleaner design keeps
a kernel-side ring buffer, has the interrupt write only there, and copies into user memory later,
when that process is the one actually running. That removes the \reg{CR3} manipulation from
interrupt context entirely.
\end{concept}

\begin{concept}{Where the return value comes from}
Look at \cod{traversal\_process->regs.eax = count;}. The count is written into the \emph{saved}
register state of the sleeping process. Later, when the scheduler picks it, \cod{load\_context}
copies that into the \cod{registers\_t} on the stack, and \cod{iret} restores it into the real
\reg{EAX}.

The user program called \cod{read()} and sees a number come back in \reg{EAX} --- with no idea that
between the call and the return, seconds passed, other programs ran, and an interrupt filled its
buffer. That illusion is the whole point of an operating system.
\end{concept}

\begin{danger}{The buffer is never NUL-terminated}
The kernel fills \cod{count} bytes and returns \cod{count}. It never writes a \cod{'\textbackslash
0'}. That is correct behaviour --- it matches Unix \cod{read()} --- but the shell then does:
\begin{lstlisting}[style=c]
char input[50];
read(0, input, 50);
input[strlen(input) - 1] = '\0';
\end{lstlisting}
\cod{strlen} on a buffer that was never terminated, and whose stack memory was never initialised,
walks until it happens to find a zero. It works in practice because a freshly allocated stack page
is zeroed, but it is reading uninitialised memory by luck. The correct form uses the return value:
\begin{lstlisting}[style=c]
int n = read(0, input, sizeof(input) - 1);
if (n > 0 && input[n-1] == '\n') n--;
input[n] = '\0';
\end{lstlisting}
Part XIII lists this among the fixes worth making.
\end{danger}

\section{Special keys}

\filetag{lytlnybl/kernel/drivers/keyboard.c}
\begin{lstlisting}[style=c]
switch (c[0]) {
    case 27:                 /* Escape */
        vga_text_clear(&terminal);
        break;
    case '\t':
        vga_text_write(&terminal, "    ");
        break;
    case '\b':
        vga_text_backspace(&terminal);
        /* deliberate fall-through */
    case '\n':
    default:
        /* ... echo and deliver to the blocked process ... */
        break;
}
\end{lstlisting}

The missing \cod{break} after backspace is intentional: backspace must erase on screen \emph{and}
still travel down to the process so it can shorten its buffer. The code below guards with
\cod{if (c[0] != '\textbackslash b')} to avoid echoing it twice.

It is correct, but it is exactly the kind of fall-through that a compiler warns about
(\cod{-Wimplicit-fallthrough}) and that a reader assumes is a bug. A one-line comment saying
``intentional'' would save the next person several minutes.

\begin{tryit}
Boot the system and press Escape at the shell prompt. The screen clears. That path goes: key
$\rightarrow$ PIC $\rightarrow$ IRQ1 $\rightarrow$ \cod{irq\_handler} $\rightarrow$
\cod{keyboard\_handler} $\rightarrow$ \cod{vga\_text\_clear}. Five subsystems for one keystroke, and
all of them are yours.
\end{tryit}

\begin{recipe}{reading a line properly}{ch24-keyboard-input}
A user program that reads a line, trims the newline, terminates it, and prints it back. Four lines
of code, and they are the four the shell gets wrong.

It reads \cod{sizeof(buf) - 1} to leave room for the terminator, uses the \emph{return value} rather
than \cod{strlen} on an unterminated buffer, and only then writes the \cod{'\textbackslash 0'}.
Compare it with the shell's version and the difference is the whole bug.
\end{recipe}

\begin{summary}
The system has a heartbeat and it has ears. It can react to time and to the user. What it still
cannot do is manage memory properly --- and without that, there are no real processes and no
isolation. That is the next part, and the most important one in the course.
\end{summary}
